% Preâmbulo copiado do template.tex do projeto do Overleaf (projeto-final-pos-rag-futebol), lido em 10/10/2026.
% Acréscimos ao template: booktabs, array, caption e url (tabelas e números gerados).
% Acréscimos pedidos pelo autor: babel em português e \raggedcolumns (o template não os tem).
\documentclass[a4paper,12pt]{article}
\usepackage[top=2cm, bottom=2cm, left=1.5cm, right=1.5cm]{geometry}
\usepackage{multicol}
\usepackage{graphicx}
\usepackage{xcolor}
\usepackage{fancyhdr}
\usepackage{tikz}
\usepackage{pgfplots}
\usepackage{pgfplotstable}
\usepackage{subcaption}
\pgfplotsset{compat=1.18}
\usetikzlibrary{shapes.geometric, positioning, arrows.meta}
\usepackage{authblk}
\usepackage{abstract}
\usepackage{float}
\usepackage{booktabs}
\usepackage{array}
\usepackage{caption}
\usepackage{url}
\usepackage[brazilian]{babel}

\setlength{\absleftindent}{0pt}
\setlength{\absrightindent}{0pt}

\renewcommand{\refname}{Referências}

\graphicspath{{generated/figures/}}
\newcommand{\modelName}{Claude Haiku 5.5} % modelo avaliado: trocar aqui se mudar
\newcommand{\embedModel}{mistral-embed} % modelo de embeddings da Busca vetorial: trocar aqui se mudar
\newcommand{\repoURL}{\url{https://github.com/joaowinderfeldbussolotto/football-tactics-knowledge-graph}} % repositório do código: trocar aqui se mudar
% Generated by scripts/paper_numbers.py from data/benchmark/final_results.jsonl. Do not edit by hand.
% Every number in the paper text is one of these macros.
\newcommand{\abstEventsInPrompt}{33\%}
\newcommand{\abstGraphTools}{2\%}
\newcommand{\abstHaikuEventsInPrompt}{33\%}
\newcommand{\abstHaikuGraphTools}{2\%}
\newcommand{\abstHaikuNoContext}{90\%}
\newcommand{\abstHaikuStatsInPrompt}{90\%}
\newcommand{\abstHaikuTextToCypher}{6\%}
\newcommand{\abstHaikuVector}{62\%}
\newcommand{\abstKimiEventsInPrompt}{0\%}
\newcommand{\abstKimiGraphTools}{4\%}
\newcommand{\abstKimiNoContext}{90\%}
\newcommand{\abstKimiStatsInPrompt}{85\%}
\newcommand{\abstKimiTextToCypher}{0\%}
\newcommand{\abstKimiVector}{33\%}
\newcommand{\abstNoContext}{90\%}
\newcommand{\abstQwenEventsInPrompt}{12\%}
\newcommand{\abstQwenGraphTools}{0\%}
\newcommand{\abstQwenNoContext}{90\%}
\newcommand{\abstQwenStatsInPrompt}{83\%}
\newcommand{\abstQwenTextToCypher}{2\%}
\newcommand{\abstQwenVector}{38\%}
\newcommand{\abstSonnetEventsInPrompt}{2\%}
\newcommand{\abstSonnetGraphTools}{2\%}
\newcommand{\abstSonnetNoContext}{90\%}
\newcommand{\abstSonnetStatsInPrompt}{83\%}
\newcommand{\abstSonnetTextToCypher}{0\%}
\newcommand{\abstSonnetVector}{58\%}
\newcommand{\abstStatsInPrompt}{90\%}
\newcommand{\abstTextToCypher}{6\%}
\newcommand{\abstVector}{62\%}
\newcommand{\accAllEventsInPrompt}{50\%}
\newcommand{\accAllGraphTools}{87\%}
\newcommand{\accAllHaikuEventsInPrompt}{50\%}
\newcommand{\accAllHaikuGraphTools}{87\%}
\newcommand{\accAllHaikuNoContext}{10\%}
\newcommand{\accAllHaikuStatsInPrompt}{10\%}
\newcommand{\accAllHaikuTextToCypher}{71\%}
\newcommand{\accAllHaikuVector}{19\%}
\newcommand{\accAllKimiEventsInPrompt}{87\%}
\newcommand{\accAllKimiGraphTools}{90\%}
\newcommand{\accAllKimiNoContext}{10\%}
\newcommand{\accAllKimiStatsInPrompt}{15\%}
\newcommand{\accAllKimiTextToCypher}{92\%}
\newcommand{\accAllKimiVector}{25\%}
\newcommand{\accAllNoContext}{10\%}
\newcommand{\accAllQwenEventsInPrompt}{44\%}
\newcommand{\accAllQwenGraphTools}{85\%}
\newcommand{\accAllQwenNoContext}{10\%}
\newcommand{\accAllQwenStatsInPrompt}{13\%}
\newcommand{\accAllQwenTextToCypher}{60\%}
\newcommand{\accAllQwenVector}{17\%}
\newcommand{\accAllSonnetEventsInPrompt}{75\%}
\newcommand{\accAllSonnetGraphTools}{90\%}
\newcommand{\accAllSonnetNoContext}{10\%}
\newcommand{\accAllSonnetStatsInPrompt}{17\%}
\newcommand{\accAllSonnetTextToCypher}{88\%}
\newcommand{\accAllSonnetVector}{23\%}
\newcommand{\accAllStatsInPrompt}{10\%}
\newcommand{\accAllTextToCypher}{71\%}
\newcommand{\accAllVector}{19\%}
\newcommand{\accControlEventsInPrompt}{100\%}
\newcommand{\accControlGraphTools}{80\%}
\newcommand{\accControlHaikuEventsInPrompt}{100\%}
\newcommand{\accControlHaikuGraphTools}{80\%}
\newcommand{\accControlHaikuNoContext}{100\%}
\newcommand{\accControlHaikuStatsInPrompt}{100\%}
\newcommand{\accControlHaikuTextToCypher}{80\%}
\newcommand{\accControlHaikuVector}{100\%}
\newcommand{\accControlKimiEventsInPrompt}{100\%}
\newcommand{\accControlKimiGraphTools}{100\%}
\newcommand{\accControlKimiNoContext}{100\%}
\newcommand{\accControlKimiStatsInPrompt}{100\%}
\newcommand{\accControlKimiTextToCypher}{80\%}
\newcommand{\accControlKimiVector}{100\%}
\newcommand{\accControlNoContext}{100\%}
\newcommand{\accControlQwenEventsInPrompt}{80\%}
\newcommand{\accControlQwenGraphTools}{100\%}
\newcommand{\accControlQwenNoContext}{100\%}
\newcommand{\accControlQwenStatsInPrompt}{100\%}
\newcommand{\accControlQwenTextToCypher}{80\%}
\newcommand{\accControlQwenVector}{100\%}
\newcommand{\accControlSonnetEventsInPrompt}{100\%}
\newcommand{\accControlSonnetGraphTools}{100\%}
\newcommand{\accControlSonnetNoContext}{100\%}
\newcommand{\accControlSonnetStatsInPrompt}{100\%}
\newcommand{\accControlSonnetTextToCypher}{80\%}
\newcommand{\accControlSonnetVector}{100\%}
\newcommand{\accControlStatsInPrompt}{100\%}
\newcommand{\accControlTextToCypher}{80\%}
\newcommand{\accControlVector}{100\%}
\newcommand{\accEventsInPrompt}{45\%}
\newcommand{\accGraphTools}{87\%}
\newcommand{\accHaikuEventsInPrompt}{45\%}
\newcommand{\accHaikuGraphTools}{87\%}
\newcommand{\accHaikuNoContext}{0\%}
\newcommand{\accHaikuStatsInPrompt}{0\%}
\newcommand{\accHaikuTextToCypher}{70\%}
\newcommand{\accHaikuVector}{11\%}
\newcommand{\accKimiEventsInPrompt}{85\%}
\newcommand{\accKimiGraphTools}{89\%}
\newcommand{\accKimiNoContext}{0\%}
\newcommand{\accKimiStatsInPrompt}{6\%}
\newcommand{\accKimiTextToCypher}{94\%}
\newcommand{\accKimiVector}{17\%}
\newcommand{\accMomentsEventsInPrompt}{17\%}
\newcommand{\accMomentsGraphTools}{100\%}
\newcommand{\accMomentsNoContext}{0\%}
\newcommand{\accMomentsStatsInPrompt}{0\%}
\newcommand{\accMomentsTextToCypher}{100\%}
\newcommand{\accMomentsVector}{0\%}
\newcommand{\accNoContext}{0\%}
\newcommand{\accPlaysEventsInPrompt}{30\%}
\newcommand{\accPlaysGraphTools}{80\%}
\newcommand{\accPlaysNoContext}{0\%}
\newcommand{\accPlaysStatsInPrompt}{0\%}
\newcommand{\accPlaysTextToCypher}{40\%}
\newcommand{\accPlaysVector}{20\%}
\newcommand{\accQwenEventsInPrompt}{40\%}
\newcommand{\accQwenGraphTools}{83\%}
\newcommand{\accQwenNoContext}{0\%}
\newcommand{\accQwenStatsInPrompt}{4\%}
\newcommand{\accQwenTextToCypher}{57\%}
\newcommand{\accQwenVector}{9\%}
\newcommand{\accRelationsEventsInPrompt}{40\%}
\newcommand{\accRelationsGraphTools}{95\%}
\newcommand{\accRelationsNoContext}{0\%}
\newcommand{\accRelationsStatsInPrompt}{0\%}
\newcommand{\accRelationsTextToCypher}{75\%}
\newcommand{\accRelationsVector}{5\%}
\newcommand{\accSearchEventsInPrompt}{82\%}
\newcommand{\accSearchGraphTools}{73\%}
\newcommand{\accSearchNoContext}{0\%}
\newcommand{\accSearchStatsInPrompt}{0\%}
\newcommand{\accSearchTextToCypher}{73\%}
\newcommand{\accSearchVector}{18\%}
\newcommand{\accSonnetEventsInPrompt}{72\%}
\newcommand{\accSonnetGraphTools}{89\%}
\newcommand{\accSonnetNoContext}{0\%}
\newcommand{\accSonnetStatsInPrompt}{9\%}
\newcommand{\accSonnetTextToCypher}{89\%}
\newcommand{\accSonnetVector}{15\%}
\newcommand{\accStatsInPrompt}{0\%}
\newcommand{\accTextToCypher}{70\%}
\newcommand{\accVector}{11\%}
\newcommand{\armEventsInPrompt}{Contexto longo}
\newcommand{\armGraphTools}{Ferramentas sobre o grafo}
\newcommand{\armNoContext}{Sem dados}
\newcommand{\armShortEventsInPrompt}{Contexto}
\newcommand{\armShortGraphTools}{Ferramentas}
\newcommand{\armShortNoContext}{Sem dados}
\newcommand{\armShortStatsInPrompt}{Estatísticas}
\newcommand{\armShortTextToCypher}{Consulta}
\newcommand{\armShortVector}{Vetorial}
\newcommand{\armStatsInPrompt}{Estatísticas resumidas}
\newcommand{\armTextToCypher}{Consulta escrita pelo modelo}
\newcommand{\armVector}{Busca vetorial}
\newcommand{\callsEventsInPrompt}{0,0}
\newcommand{\callsGraphTools}{3,3}
\newcommand{\callsHaikuEventsInPrompt}{0,0}
\newcommand{\callsHaikuGraphTools}{3,3}
\newcommand{\callsHaikuNoContext}{0,0}
\newcommand{\callsHaikuStatsInPrompt}{0,0}
\newcommand{\callsHaikuTextToCypher}{4,8}
\newcommand{\callsHaikuVector}{0,0}
\newcommand{\callsKimiEventsInPrompt}{0,0}
\newcommand{\callsKimiGraphTools}{3,8}
\newcommand{\callsKimiNoContext}{0,0}
\newcommand{\callsKimiStatsInPrompt}{0,0}
\newcommand{\callsKimiTextToCypher}{3,0}
\newcommand{\callsKimiVector}{0,0}
\newcommand{\callsNoContext}{0,0}
\newcommand{\callsQwenEventsInPrompt}{0,0}
\newcommand{\callsQwenGraphTools}{6,2}
\newcommand{\callsQwenNoContext}{0,0}
\newcommand{\callsQwenStatsInPrompt}{0,0}
\newcommand{\callsQwenTextToCypher}{10,5}
\newcommand{\callsQwenVector}{0,0}
\newcommand{\callsSonnetEventsInPrompt}{0,0}
\newcommand{\callsSonnetGraphTools}{1,8}
\newcommand{\callsSonnetNoContext}{0,0}
\newcommand{\callsSonnetStatsInPrompt}{0,0}
\newcommand{\callsSonnetTextToCypher}{2,0}
\newcommand{\callsSonnetVector}{0,0}
\newcommand{\callsStatsInPrompt}{0,0}
\newcommand{\callsTextToCypher}{4,8}
\newcommand{\callsVector}{0,0}
\newcommand{\ciHighEventsInPrompt}{59}
\newcommand{\ciHighGraphTools}{94}
\newcommand{\ciHighHaikuEventsInPrompt}{59}
\newcommand{\ciHighHaikuGraphTools}{94}
\newcommand{\ciHighHaikuNoContext}{8}
\newcommand{\ciHighHaikuStatsInPrompt}{8}
\newcommand{\ciHighHaikuTextToCypher}{81}
\newcommand{\ciHighHaikuVector}{23}
\newcommand{\ciHighKimiEventsInPrompt}{93}
\newcommand{\ciHighKimiGraphTools}{95}
\newcommand{\ciHighKimiNoContext}{8}
\newcommand{\ciHighKimiStatsInPrompt}{17}
\newcommand{\ciHighKimiTextToCypher}{98}
\newcommand{\ciHighKimiVector}{30}
\newcommand{\ciHighNoContext}{8}
\newcommand{\ciHighQwenEventsInPrompt}{55}
\newcommand{\ciHighQwenGraphTools}{91}
\newcommand{\ciHighQwenNoContext}{8}
\newcommand{\ciHighQwenStatsInPrompt}{14}
\newcommand{\ciHighQwenTextToCypher}{70}
\newcommand{\ciHighQwenVector}{20}
\newcommand{\ciHighSonnetEventsInPrompt}{83}
\newcommand{\ciHighSonnetGraphTools}{95}
\newcommand{\ciHighSonnetNoContext}{8}
\newcommand{\ciHighSonnetStatsInPrompt}{20}
\newcommand{\ciHighSonnetTextToCypher}{95}
\newcommand{\ciHighSonnetVector}{28}
\newcommand{\ciHighStatsInPrompt}{8}
\newcommand{\ciHighTextToCypher}{81}
\newcommand{\ciHighVector}{23}
\newcommand{\ciLevel}{95\%}
\newcommand{\ciLowEventsInPrompt}{31}
\newcommand{\ciLowGraphTools}{75}
\newcommand{\ciLowHaikuEventsInPrompt}{31}
\newcommand{\ciLowHaikuGraphTools}{75}
\newcommand{\ciLowHaikuNoContext}{0}
\newcommand{\ciLowHaikuStatsInPrompt}{0}
\newcommand{\ciLowHaikuTextToCypher}{56}
\newcommand{\ciLowHaikuVector}{5}
\newcommand{\ciLowKimiEventsInPrompt}{72}
\newcommand{\ciLowKimiGraphTools}{77}
\newcommand{\ciLowKimiNoContext}{0}
\newcommand{\ciLowKimiStatsInPrompt}{2}
\newcommand{\ciLowKimiTextToCypher}{83}
\newcommand{\ciLowKimiVector}{9}
\newcommand{\ciLowNoContext}{0}
\newcommand{\ciLowQwenEventsInPrompt}{28}
\newcommand{\ciLowQwenGraphTools}{70}
\newcommand{\ciLowQwenNoContext}{0}
\newcommand{\ciLowQwenStatsInPrompt}{1}
\newcommand{\ciLowQwenTextToCypher}{43}
\newcommand{\ciLowQwenVector}{3}
\newcommand{\ciLowSonnetEventsInPrompt}{58}
\newcommand{\ciLowSonnetGraphTools}{77}
\newcommand{\ciLowSonnetNoContext}{0}
\newcommand{\ciLowSonnetStatsInPrompt}{3}
\newcommand{\ciLowSonnetTextToCypher}{77}
\newcommand{\ciLowSonnetVector}{7}
\newcommand{\ciLowStatsInPrompt}{0}
\newcommand{\ciLowTextToCypher}{56}
\newcommand{\ciLowVector}{5}
\newcommand{\costFinal}{1,04}
\newcommand{\costHaiku}{1,04}
\newcommand{\costKimi}{33,40}
\newcommand{\costQwen}{0,48}
\newcommand{\costSonnet}{28,01}
\newcommand{\fmtEventsInPrompt}{0\%}
\newcommand{\fmtGraphTools}{0\%}
\newcommand{\fmtHaikuEventsInPrompt}{0\%}
\newcommand{\fmtHaikuGraphTools}{0\%}
\newcommand{\fmtHaikuNoContext}{0\%}
\newcommand{\fmtHaikuStatsInPrompt}{0\%}
\newcommand{\fmtHaikuTextToCypher}{0\%}
\newcommand{\fmtHaikuVector}{0\%}
\newcommand{\fmtKimiEventsInPrompt}{0\%}
\newcommand{\fmtKimiGraphTools}{4\%}
\newcommand{\fmtKimiNoContext}{0\%}
\newcommand{\fmtKimiStatsInPrompt}{0\%}
\newcommand{\fmtKimiTextToCypher}{0\%}
\newcommand{\fmtKimiVector}{0\%}
\newcommand{\fmtNoContext}{0\%}
\newcommand{\fmtQwenEventsInPrompt}{0\%}
\newcommand{\fmtQwenGraphTools}{6\%}
\newcommand{\fmtQwenNoContext}{0\%}
\newcommand{\fmtQwenStatsInPrompt}{2\%}
\newcommand{\fmtQwenTextToCypher}{15\%}
\newcommand{\fmtQwenVector}{2\%}
\newcommand{\fmtSonnetEventsInPrompt}{4\%}
\newcommand{\fmtSonnetGraphTools}{0\%}
\newcommand{\fmtSonnetNoContext}{0\%}
\newcommand{\fmtSonnetStatsInPrompt}{0\%}
\newcommand{\fmtSonnetTextToCypher}{0\%}
\newcommand{\fmtSonnetVector}{0\%}
\newcommand{\fmtStatsInPrompt}{0\%}
\newcommand{\fmtTextToCypher}{0\%}
\newcommand{\fmtVector}{0\%}
\newcommand{\hallEventsInPrompt}{17\%}
\newcommand{\hallGraphTools}{12\%}
\newcommand{\hallHaikuEventsInPrompt}{17\%}
\newcommand{\hallHaikuGraphTools}{12\%}
\newcommand{\hallHaikuNoContext}{0\%}
\newcommand{\hallHaikuStatsInPrompt}{0\%}
\newcommand{\hallHaikuTextToCypher}{23\%}
\newcommand{\hallHaikuVector}{19\%}
\newcommand{\hallKimiEventsInPrompt}{13\%}
\newcommand{\hallKimiGraphTools}{2\%}
\newcommand{\hallKimiNoContext}{0\%}
\newcommand{\hallKimiStatsInPrompt}{0\%}
\newcommand{\hallKimiTextToCypher}{8\%}
\newcommand{\hallKimiVector}{42\%}
\newcommand{\hallNoContext}{0\%}
\newcommand{\hallQwenEventsInPrompt}{44\%}
\newcommand{\hallQwenGraphTools}{10\%}
\newcommand{\hallQwenNoContext}{0\%}
\newcommand{\hallQwenStatsInPrompt}{2\%}
\newcommand{\hallQwenTextToCypher}{23\%}
\newcommand{\hallQwenVector}{42\%}
\newcommand{\hallSonnetEventsInPrompt}{19\%}
\newcommand{\hallSonnetGraphTools}{8\%}
\newcommand{\hallSonnetNoContext}{0\%}
\newcommand{\hallSonnetStatsInPrompt}{0\%}
\newcommand{\hallSonnetTextToCypher}{12\%}
\newcommand{\hallSonnetVector}{19\%}
\newcommand{\hallStatsInPrompt}{0\%}
\newcommand{\hallTextToCypher}{23\%}
\newcommand{\hallVector}{19\%}
\newcommand{\kAllEventsInPrompt}{26}
\newcommand{\kAllGraphTools}{45}
\newcommand{\kAllNoContext}{5}
\newcommand{\kAllStatsInPrompt}{5}
\newcommand{\kAllTextToCypher}{37}
\newcommand{\kAllVector}{10}
\newcommand{\kControlEventsInPrompt}{5}
\newcommand{\kControlGraphTools}{4}
\newcommand{\kControlNoContext}{5}
\newcommand{\kControlStatsInPrompt}{5}
\newcommand{\kControlTextToCypher}{4}
\newcommand{\kControlVector}{5}
\newcommand{\kEventsInPrompt}{21}
\newcommand{\kGraphTools}{41}
\newcommand{\kHaikuEventsInPrompt}{21}
\newcommand{\kHaikuGraphTools}{41}
\newcommand{\kHaikuNoContext}{0}
\newcommand{\kHaikuStatsInPrompt}{0}
\newcommand{\kHaikuTextToCypher}{33}
\newcommand{\kHaikuVector}{5}
\newcommand{\kKimiEventsInPrompt}{40}
\newcommand{\kKimiGraphTools}{42}
\newcommand{\kKimiNoContext}{0}
\newcommand{\kKimiStatsInPrompt}{3}
\newcommand{\kKimiTextToCypher}{44}
\newcommand{\kKimiVector}{8}
\newcommand{\kMomentsEventsInPrompt}{1}
\newcommand{\kMomentsGraphTools}{6}
\newcommand{\kMomentsNoContext}{0}
\newcommand{\kMomentsStatsInPrompt}{0}
\newcommand{\kMomentsTextToCypher}{6}
\newcommand{\kMomentsVector}{0}
\newcommand{\kNoContext}{0}
\newcommand{\kPlaysEventsInPrompt}{3}
\newcommand{\kPlaysGraphTools}{8}
\newcommand{\kPlaysNoContext}{0}
\newcommand{\kPlaysStatsInPrompt}{0}
\newcommand{\kPlaysTextToCypher}{4}
\newcommand{\kPlaysVector}{2}
\newcommand{\kQwenEventsInPrompt}{19}
\newcommand{\kQwenGraphTools}{39}
\newcommand{\kQwenNoContext}{0}
\newcommand{\kQwenStatsInPrompt}{2}
\newcommand{\kQwenTextToCypher}{27}
\newcommand{\kQwenVector}{4}
\newcommand{\kRelationsEventsInPrompt}{8}
\newcommand{\kRelationsGraphTools}{19}
\newcommand{\kRelationsNoContext}{0}
\newcommand{\kRelationsStatsInPrompt}{0}
\newcommand{\kRelationsTextToCypher}{15}
\newcommand{\kRelationsVector}{1}
\newcommand{\kSearchEventsInPrompt}{9}
\newcommand{\kSearchGraphTools}{8}
\newcommand{\kSearchNoContext}{0}
\newcommand{\kSearchStatsInPrompt}{0}
\newcommand{\kSearchTextToCypher}{8}
\newcommand{\kSearchVector}{2}
\newcommand{\kSonnetEventsInPrompt}{34}
\newcommand{\kSonnetGraphTools}{42}
\newcommand{\kSonnetNoContext}{0}
\newcommand{\kSonnetStatsInPrompt}{4}
\newcommand{\kSonnetTextToCypher}{42}
\newcommand{\kSonnetVector}{7}
\newcommand{\kStatsInPrompt}{0}
\newcommand{\kTextToCypher}{33}
\newcommand{\kVector}{5}
\newcommand{\latEventsInPrompt}{3,9}
\newcommand{\latGraphTools}{7,8}
\newcommand{\latHaikuEventsInPrompt}{3,9}
\newcommand{\latHaikuGraphTools}{7,8}
\newcommand{\latHaikuNoContext}{1,8}
\newcommand{\latHaikuStatsInPrompt}{1,7}
\newcommand{\latHaikuTextToCypher}{10,1}
\newcommand{\latHaikuVector}{2,3}
\newcommand{\latKimiEventsInPrompt}{71,0}
\newcommand{\latKimiGraphTools}{81,1}
\newcommand{\latKimiNoContext}{4,4}
\newcommand{\latKimiStatsInPrompt}{6,2}
\newcommand{\latKimiTextToCypher}{38,2}
\newcommand{\latKimiVector}{25,3}
\newcommand{\latNoContext}{1,8}
\newcommand{\latQwenEventsInPrompt}{28,3}
\newcommand{\latQwenGraphTools}{31,7}
\newcommand{\latQwenNoContext}{4,1}
\newcommand{\latQwenStatsInPrompt}{4,7}
\newcommand{\latQwenTextToCypher}{49,2}
\newcommand{\latQwenVector}{6,2}
\newcommand{\latSonnetEventsInPrompt}{51,5}
\newcommand{\latSonnetGraphTools}{18,7}
\newcommand{\latSonnetNoContext}{4,5}
\newcommand{\latSonnetStatsInPrompt}{5,1}
\newcommand{\latSonnetTextToCypher}{16,0}
\newcommand{\latSonnetVector}{14,2}
\newcommand{\latStatsInPrompt}{1,7}
\newcommand{\latTextToCypher}{10,1}
\newcommand{\latVector}{2,3}
\newcommand{\mcGraphToolsEventsInPromptFirst}{23}
\newcommand{\mcGraphToolsEventsInPromptP}{menor que 0,001}
\newcommand{\mcGraphToolsEventsInPromptPHolm}{menor que 0,001}
\newcommand{\mcGraphToolsEventsInPromptSecond}{3}
\newcommand{\mcGraphToolsTextToCypherFirst}{11}
\newcommand{\mcGraphToolsTextToCypherP}{0,057}
\newcommand{\mcGraphToolsTextToCypherPHolm}{0,057}
\newcommand{\mcGraphToolsTextToCypherSecond}{3}
\newcommand{\mcHaikuGraphToolsEventsInPromptFirst}{23}
\newcommand{\mcHaikuGraphToolsEventsInPromptP}{menor que 0,001}
\newcommand{\mcHaikuGraphToolsEventsInPromptPHolm}{menor que 0,001}
\newcommand{\mcHaikuGraphToolsEventsInPromptSecond}{3}
\newcommand{\mcHaikuGraphToolsTextToCypherFirst}{11}
\newcommand{\mcHaikuGraphToolsTextToCypherP}{0,057}
\newcommand{\mcHaikuGraphToolsTextToCypherPHolm}{0,057}
\newcommand{\mcHaikuGraphToolsTextToCypherSecond}{3}
\newcommand{\mcHaikuTextToCypherEventsInPromptFirst}{17}
\newcommand{\mcHaikuTextToCypherEventsInPromptP}{0,017}
\newcommand{\mcHaikuTextToCypherEventsInPromptPHolm}{0,034}
\newcommand{\mcHaikuTextToCypherEventsInPromptSecond}{5}
\newcommand{\mcKimiGraphToolsEventsInPromptFirst}{6}
\newcommand{\mcKimiGraphToolsEventsInPromptP}{0,754}
\newcommand{\mcKimiGraphToolsEventsInPromptPHolm}{1,000}
\newcommand{\mcKimiGraphToolsEventsInPromptSecond}{4}
\newcommand{\mcKimiGraphToolsTextToCypherFirst}{2}
\newcommand{\mcKimiGraphToolsTextToCypherP}{0,688}
\newcommand{\mcKimiGraphToolsTextToCypherPHolm}{1,000}
\newcommand{\mcKimiGraphToolsTextToCypherSecond}{4}
\newcommand{\mcKimiTextToCypherEventsInPromptFirst}{6}
\newcommand{\mcKimiTextToCypherEventsInPromptP}{0,289}
\newcommand{\mcKimiTextToCypherEventsInPromptPHolm}{0,867}
\newcommand{\mcKimiTextToCypherEventsInPromptSecond}{2}
\newcommand{\mcQwenGraphToolsEventsInPromptFirst}{23}
\newcommand{\mcQwenGraphToolsEventsInPromptP}{menor que 0,001}
\newcommand{\mcQwenGraphToolsEventsInPromptPHolm}{menor que 0,001}
\newcommand{\mcQwenGraphToolsEventsInPromptSecond}{3}
\newcommand{\mcQwenGraphToolsTextToCypherFirst}{15}
\newcommand{\mcQwenGraphToolsTextToCypherP}{0,008}
\newcommand{\mcQwenGraphToolsTextToCypherPHolm}{0,015}
\newcommand{\mcQwenGraphToolsTextToCypherSecond}{3}
\newcommand{\mcQwenTextToCypherEventsInPromptFirst}{12}
\newcommand{\mcQwenTextToCypherEventsInPromptP}{0,077}
\newcommand{\mcQwenTextToCypherEventsInPromptPHolm}{0,077}
\newcommand{\mcQwenTextToCypherEventsInPromptSecond}{4}
\newcommand{\mcSonnetGraphToolsEventsInPromptFirst}{12}
\newcommand{\mcSonnetGraphToolsEventsInPromptP}{0,077}
\newcommand{\mcSonnetGraphToolsEventsInPromptPHolm}{0,154}
\newcommand{\mcSonnetGraphToolsEventsInPromptSecond}{4}
\newcommand{\mcSonnetGraphToolsTextToCypherFirst}{5}
\newcommand{\mcSonnetGraphToolsTextToCypherP}{1,000}
\newcommand{\mcSonnetGraphToolsTextToCypherPHolm}{1,000}
\newcommand{\mcSonnetGraphToolsTextToCypherSecond}{5}
\newcommand{\mcSonnetTextToCypherEventsInPromptFirst}{8}
\newcommand{\mcSonnetTextToCypherEventsInPromptP}{0,008}
\newcommand{\mcSonnetTextToCypherEventsInPromptPHolm}{0,023}
\newcommand{\mcSonnetTextToCypherEventsInPromptSecond}{0}
\newcommand{\mcTextToCypherEventsInPromptFirst}{17}
\newcommand{\mcTextToCypherEventsInPromptP}{0,017}
\newcommand{\mcTextToCypherEventsInPromptPHolm}{0,034}
\newcommand{\mcTextToCypherEventsInPromptSecond}{5}
\newcommand{\modelHaiku}{Claude Haiku 5.5}
\newcommand{\modelKimi}{Kimi K3}
\newcommand{\modelList}{Claude Haiku 5.5, Qwen3.8-flash, Kimi K3 e Claude Sonnet 4.6}
\newcommand{\modelQwen}{Qwen3.8-flash}
\newcommand{\modelSonnet}{Claude Sonnet 4.6}
\newcommand{\nActions}{2.585}
\newcommand{\nAnswerable}{47}
\newcommand{\nArms}{6}
\newcommand{\nCheckNoData}{5}
\newcommand{\nCheckPlayer}{35}
\newcommand{\nCheckSet}{10}
\newcommand{\nCheckValue}{2}
\newcommand{\nControl}{5}
\newcommand{\nEvents}{4.407}
\newcommand{\nGroupControl}{5}
\newcommand{\nGroupMoments}{6}
\newcommand{\nGroupPlays}{10}
\newcommand{\nGroupRelations}{20}
\newcommand{\nGroupSearch}{11}
\newcommand{\nMaxCalls}{12}
\newcommand{\nMaxRows}{50}
\newcommand{\nModels}{4}
\newcommand{\nPhases}{537}
\newcommand{\nPlayers}{34}
\newcommand{\nQuestions}{52}
\newcommand{\nRepeats}{1}
\newcommand{\nRuns}{312}
\newcommand{\nTools}{8}
\newcommand{\nTopK}{30}
\newcommand{\nZones}{96}
\newcommand{\preFixKGraphTools}{39}
\newcommand{\preFixNGraphTools}{52}
\newcommand{\repEventsInPrompt}{1}
\newcommand{\repGraphTools}{1}
\newcommand{\repNoContext}{1}
\newcommand{\repStatsInPrompt}{1}
\newcommand{\repTextToCypher}{1}
\newcommand{\repVector}{1}
\newcommand{\tokInEventsInPrompt}{115.735}
\newcommand{\tokInGraphTools}{76.118}
\newcommand{\tokInHaikuEventsInPrompt}{115.735}
\newcommand{\tokInHaikuGraphTools}{76.118}
\newcommand{\tokInHaikuNoContext}{880}
\newcommand{\tokInHaikuStatsInPrompt}{4.968}
\newcommand{\tokInHaikuTextToCypher}{50.983}
\newcommand{\tokInHaikuVector}{2.623}
\newcommand{\tokInKimiEventsInPrompt}{89.129}
\newcommand{\tokInKimiGraphTools}{43.488}
\newcommand{\tokInKimiNoContext}{458}
\newcommand{\tokInKimiStatsInPrompt}{4.423}
\newcommand{\tokInKimiTextToCypher}{23.524}
\newcommand{\tokInKimiVector}{2.207}
\newcommand{\tokInNoContext}{880}
\newcommand{\tokInQwenEventsInPrompt}{98.311}
\newcommand{\tokInQwenGraphTools}{93.773}
\newcommand{\tokInQwenNoContext}{673}
\newcommand{\tokInQwenStatsInPrompt}{4.904}
\newcommand{\tokInQwenTextToCypher}{85.206}
\newcommand{\tokInQwenVector}{2.080}
\newcommand{\tokInSonnetEventsInPrompt}{90.128}
\newcommand{\tokInSonnetGraphTools}{35.059}
\newcommand{\tokInSonnetNoContext}{501}
\newcommand{\tokInSonnetStatsInPrompt}{4.052}
\newcommand{\tokInSonnetTextToCypher}{19.145}
\newcommand{\tokInSonnetVector}{1.916}
\newcommand{\tokInStatsInPrompt}{4.968}
\newcommand{\tokInTextToCypher}{50.983}
\newcommand{\tokInVector}{2.623}
\newcommand{\tokOutEventsInPrompt}{615}
\newcommand{\tokOutGraphTools}{574}
\newcommand{\tokOutNoContext}{180}
\newcommand{\tokOutStatsInPrompt}{192}
\newcommand{\tokOutTextToCypher}{1.198}
\newcommand{\tokOutVector}{338}
\newcommand{\tokTableApprox}{cerca de 100 mil tokens}

\newcommand{\gentable}[1]{\input{generated/tables/#1}} % no Overleaf: \input{tab-#1}

\title{Seis formas de recuperação de dados de eventos de futebol para LLMs} % título provisório
\author[1]{Autor A}
\author[1]{Autor B}
\author[1]{Autor C}

\affil[1]{Pós-graduação em Inteligência Artificial Aplicada, UniSENAI}

\date{}

\renewcommand\Authand{ e }
\renewcommand\Authands{ e }

\begin{document}

\maketitle

\begin{multicols*}{2}
\raggedcolumns % sem isto o LaTeX estica o espaço entre títulos para alinhar o fim das colunas

% Resumo: ambiente abstract e palavras-chave, como no template da disciplina (src/abstract.tex no Overleaf).
% Resultado só qualitativo, sem números (pedido do autor); termina no último resultado, sem frase de implicação.

\begin{abstract}
Modelos de linguagem permitem fazer perguntas em linguagem natural sobre dados de eventos de futebol, mas os dados
da partida precisam chegar ao contexto do modelo. Este trabalho transforma os eventos da final da Copa do Mundo de
2022 em uma tabela de ações e em um grafo, no qual cada ação liga um jogador à partida, e compara \nArms{}
abordagens de aprendizado em contexto para dar esses dados ao modelo: sem dados, a tabela de ações inteira no
prompt, estatísticas resumidas, busca vetorial, consulta escrita pelo modelo na linguagem do grafo (Cypher) e
ferramentas com consultas parametrizadas sobre o grafo. A avaliação usou \nQuestions{} perguntas sobre a partida,
em \nModels{} modelos, dois menores e dois maiores. As ferramentas sobre o grafo foram a única abordagem com acerto
alto em todos os modelos. A consulta escrita pelo modelo só alcançou esse nível nos modelos maiores, com raciocínio estendido. Nesses modelos, a tabela inteira no prompt também chegou perto, mas gastou mais tokens de saída e mais tempo por pergunta. A busca vetorial e as estatísticas resumidas acertaram pouco em todos os modelos.
\end{abstract}

\noindent\textbf{Palavras-chave:} modelos de linguagem, aprendizado em contexto, recuperação de dados, grafos,
dados de eventos de futebol.

% Seção: Introdução  (orçamento: 1 página)
%
% Curta e sem detalhe: o como fica na Metodologia e os números nos Resultados. Sem hipóteses formais (ALERTAS,
% item 2) e sem parágrafo de organização do texto.

\section{Introdução}\label{sec:introdução}

Os dados de eventos de futebol registram cada ação de uma partida, com o jogador, o momento e a posição no
campo, e há conjuntos desse tipo abertos ao público~\cite{pappalardo2019public}. Muitas perguntas de análise
sobre esses dados exigem ligar eventos entre si: quem mais passou para quem, em que ordem a bola circulou em uma
jogada, o que mudou depois de uma substituição.

Um LLM permite fazer essas perguntas em linguagem natural, mas os dados da partida precisam chegar ao seu
contexto. Isso pode ser feito de várias formas: com a tabela inteira no prompt, com as linhas mais parecidas com
a pergunta, com estatísticas resumidas, com uma consulta escrita pelo próprio modelo ou com ferramentas que já
trazem a consulta pronta. Em todas, o modelo não é treinado de novo; muda o que entra no contexto e quem junta
os eventos, o modelo ou o código. No esporte, os trabalhos publicados avaliam uma dessas formas de cada vez
(Seção~\ref{sec:trabalhos}).

Este trabalho compara \nArms{} abordagens sobre a final da Copa do Mundo de 2022, com \nQuestions{} perguntas
cujo gabarito foi calculado dos dados brutos, em \nModels{} modelos de tamanhos diferentes. As perguntas de
pesquisa são qual abordagem acerta mais, como o acerto muda de um modelo para outro e quanto cada abordagem
custa em tokens e tempo. \armGraphTools{} teve acerto alto em todos os modelos. \armTextToCypher{} e
\armEventsInPrompt{} só se aproximaram dela nos modelos maiores, com mais tokens de saída e mais tempo por
pergunta.

% Contribuições: subseção da Introdução, como no template da disciplina. Só o que o trabalho entrega; os
% resultados ficam no último parágrafo da Introdução e na Discussão.

\subsection{Contribuições}\label{sec:contribuições}

\begin{itemize}
\item Uma transformação dos eventos da partida em uma tabela de ações, com colunas de contexto (zona do campo,
recebedor do passe, fase de posse e placar), e em um grafo, em que cada ação é uma aresta do jogador para a
partida (Seção~\ref{sec:dados}).
\item Um conjunto de \nQuestions{} perguntas sobre a final da Copa do Mundo de 2022, com gabarito
(Seção~\ref{sec:perguntas}).
O código, as perguntas e o gabarito são públicos (Seção~\ref{sec:setup}).
\item Uma comparação de \nArms{} abordagens de aprendizado em contexto sob as mesmas perguntas, o mesmo prompt e
a mesma saída, em \nModels{} modelos de tamanhos diferentes (Seção~\ref{sec:experimentos}).
\end{itemize}

O restante do texto está organizado da seguinte forma: a Seção~\ref{sec:trabalhos} situa o trabalho em relação
à literatura; a Seção~\ref{sec:metodologia} descreve os dados, as abordagens comparadas, as perguntas e a
conferência das respostas; a Seção~\ref{sec:experimentos} apresenta os resultados e a discussão; e a
Seção~\ref{sec:conclusão} conclui o trabalho.

% Seção: Trabalhos relacionados  (orçamento: 1 a 1,5 página)
%
% Formato combinado com o autor: quatro trabalhos com um parágrafo cada (SoccerRAG, SPORTSQL, Tzikas, LOFT) e os
% demais em uma frase, mais a tabela de posicionamento. Esta seção só posiciona: comparação com os resultados
% (por exemplo, o acerto de Tzikas contra o da Busca vetorial) fica na Discussão.
% Fora daqui: NLGraph, GraphRAG (Edge), StructGPT, Pan et al. e Pappalardo; Duch, Buldú e Sui ficam para a Discussão,
% se um resultado pedir. "Pergunta encadeada" é definida na Seção 5.3; aqui a definição curta vem com a remissão.
% Números de outros trabalhos vão por extenso (o lint só aceita números do benchmark como macro).

\section{Trabalhos relacionados}\label{sec:trabalhos}

Os primeiros conjuntos de avaliação de LLMs no esporte medem o que o modelo sabe, sem dados de uma partida. O
SportQA~\cite{xia2024sportqa} reúne mais de setenta mil perguntas de múltipla escolha, de fatos históricos a
cenários de jogo; os modelos avaliados foram bem nos fatos básicos e ficaram atrás de especialistas humanos nos
cenários. O SoccerAgent~\cite{rao2025socceragent} divide as perguntas entre vários agentes que consultam uma base
de conhecimento sobre jogadores, times, árbitros e estádios, também com perguntas de múltipla escolha.

Os trabalhos mais próximos ligam o modelo a dados de partidas. O SoccerRAG~\cite{strand2024soccerrag} guarda em
um banco SQLite as anotações de eventos do SoccerNet, como gols, cartões e faltas, os comentários transcritos das
transmissões e os dados dos jogadores. O modelo extrai e valida os nomes citados na pergunta, e um agente escreve
a consulta SQL. A avaliação usou vinte perguntas de exemplo e comparou seis configurações do próprio sistema em
dez delas, todas com o modelo escrevendo SQL.

O SPORTSQL~\cite{martinez2025sportsql} traduz perguntas sobre a liga inglesa em consultas SQL sobre um banco
montado com os dados do Fantasy Premier League e indexado no tempo. Acompanha um conjunto de mais de mil e
setecentas perguntas, cada uma com a consulta SQL e a resposta de referência, o que mostra que a conferência por
resposta exata é viável no esporte. Os dados, porém, são estatísticas de jogadores e times, e não o registro de
cada ação da partida.

Tzikas et al.~\cite{tzikas2025dynamic} usam eventos de cinco rodadas da liga alemã, com posição no campo, tempo e
jogador. Cada evento vira uma frase montada a partir de um molde de texto, e a resposta vem de recuperação
aumentada por geração, com reordenação dos trechos recuperados. As trinta perguntas são de três tipos (quando,
onde e quem), e a conferência é feita por um LLM, que julga se a resposta tem o mesmo sentido do gabarito. Os
autores relatam erros desse julgamento e apontam a recuperação iterativa como caminho para as perguntas que
exigem juntar informação espalhada pela partida.

Fora do esporte, a comparação entre formas de recuperação já foi estudada em dados gerais.
Li et al.~\cite{li2024ragorlongcontext} compararam recuperação e contexto longo em conjuntos públicos: com
recursos suficientes, o contexto longo teve desempenho médio maior, e a recuperação, custo bem menor. Liu
et al.~\cite{liu2024lostinthemiddle} mostraram que o desempenho cai quando a informação relevante está no meio de
um contexto longo. Han et al.~\cite{han2025ragvsgraphrag} compararam a recuperação de trechos de texto com a
recuperação ao longo de um grafo, sob um mesmo protocolo, e cada forma se saiu melhor em tarefas diferentes.
Quando os dados estão em um banco de grafos, o modelo também pode escrever a consulta, mas LLMs ainda erram ao
gerar consultas Cypher~\cite{ozsoy2025text2cypher}.

O LOFT~\cite{lee2024loft} é o que mais se aproxima da comparação feita aqui. O conjunto testa se um modelo que
recebe o corpus inteiro no contexto substitui sistemas de recuperação, de recuperação aumentada por geração e de
SQL. O contexto longo teve desempenho comparável ao dos sistemas de recuperação, mas teve dificuldade com o
raciocínio composicional das tarefas do tipo SQL, isto é, as que combinam várias partes dos dados.

\vspace{2mm}
\noindent\begin{minipage}{\columnwidth}
\centering\footnotesize\renewcommand{\arraystretch}{1.15}\setlength{\tabcolsep}{3pt}
\begin{tabular}{>{\raggedright\arraybackslash}p{0.25\columnwidth}>{\raggedright\arraybackslash}p{0.235\columnwidth}>{\raggedright\arraybackslash}p{0.21\columnwidth}>{\raggedright\arraybackslash}p{0.21\columnwidth}} % lint: ignore
\toprule
Trabalho & Dados & Formas de recuperação & Conferência \\
\midrule
SportQA \cite{xia2024sportqa} & nenhum (memória do modelo) & nenhuma & múltipla escolha \\
SoccerAgent \cite{rao2025socceragent} & base de conhecimento & uma (agentes) & múltipla escolha \\
SoccerRAG \cite{strand2024soccerrag} & anotações e comentários de partidas & uma (SQL), em seis configurações & feita pelos autores \\
SPORTSQL \cite{martinez2025sportsql} & estatísticas de jogadores e times & uma (SQL) & resposta de referência \\
Tzikas et al. \cite{tzikas2025dynamic} & eventos de partidas & uma (recuperação vetorial) & LLM como juiz \\
Li et al. \cite{li2024ragorlongcontext} & textos gerais & duas (recuperação e contexto longo) & métricas dos conjuntos \\
LOFT \cite{lee2024loft} & textos e tabelas gerais & contexto longo contra sistemas especializados & métricas dos conjuntos \\
\bottomrule
\end{tabular}
\captionof{table}{Trabalhos relacionados: dados, formas de recuperação e conferência da resposta.}\label{tab:posicionamento}
\end{minipage}\par\smallskip

A Tabela~\ref{tab:posicionamento} resume esses trabalhos. No esporte, cada trabalho avalia um sistema com uma
única forma de recuperação; fora dele, as formas de recuperação são comparadas, mas sem dados esportivos. Não foram
encontrados trabalhos que comparem formas de recuperação em perguntas encadeadas sobre eventos de futebol, isto é,
perguntas cuja resposta exige ligar eventos entre si (Seção~\ref{sec:perguntas}), com conferência por código
contra um gabarito calculado dos dados brutos.

% Seção: Metodologia  (orçamento: 2.5 página(s); passou disso por pedido do autor: duas figuras e três tabelas)
%
% Referências previstas (chaves de refs.bib): schroer2021crispdm, martinezplumed2021crispdm, shimaoka2024crispdm, decroos2019actions, brandes2001faster, yao2023react, ozsoy2025text2cypher, lewis2020rag, sainz2023nlp, ji2023survey
% Figuras e tabelas ainda a fazer: F5 (exemplo de robustez).
% O modelo é a macro \modelName, definida no preâmbulo (main.tex); trocar lá.
% Brandes (intermediação) ficou fora do texto a pedido do autor; a chave continua em refs.bib.

\section{Metodologia}\label{sec:metodologia}

O trabalho seguiu as seis fases do CRISP-DM~\cite{schroer2021crispdm}: entendimento do problema, entendimento
dos dados, preparação, modelagem, avaliação e implantação. A pergunta de partida é quem faz o cálculo quando um
modelo de linguagem responde sobre dados de eventos de futebol: o próprio modelo, uma consulta que ele escreve ou
código que já está escrito.

\subsection{Dados e transformações}

Os dados são os eventos abertos da StatsBomb~\cite{statsbomb_opendata} da final da Copa do Mundo de 2022. Os
\nEvents{} eventos da partida foram convertidos em \nActions{} ações no formato SPADL~\cite{decroos2019actions},
em uma tabela com uma linha por ação, com zona do campo, fase de posse e valor esperado de ameaça (xT). A tabela
foi carregada em um grafo Neo4j, do qual também se calcula uma súmula com totais por jogador e por time
(Figura~\ref{fig:dados}). Uma única partida foi usada para que a tabela inteira caiba no contexto do modelo: ela
ocupa \tokTableApprox{}, valor que varia com o tokenizador do modelo.

% F1: das transformações dos dados aos braços. Diagrama em TikZ (nenhum dado de resultado).
% Fica na raiz de paper/ porque o Overleaf do artigo não tem pastas.
\vspace{2mm}
\noindent\begin{minipage}{\columnwidth}
\centering\small
\begin{tikzpicture}[
  box/.style={draw, rounded corners=2pt, align=center, inner sep=4pt, text width=0.88\columnwidth},
  arr/.style={-{Stealth[length=2mm]}}]
\node[box] (raw) at (0,0) {Eventos brutos (\nEvents{} eventos)\\ lidos pelo gabarito};
\node[box] (act) at (0,-2.5) {Tabela de ações (\nActions{} ações)\\ lida por \armEventsInPrompt{} e \armVector{}};
\node[box] (grafo) at (0,-5.0) {Grafo Neo4j e súmula agregada\\ lidos por \armStatsInPrompt{}, \armTextToCypher{} e \armGraphTools{}};
\draw[arr] (raw) -- node[right, xshift=1mm] {SPADL, zonas, posse e xT} (act);
\draw[arr] (act) -- node[right, xshift=1mm] {carga no banco} (grafo);
\end{tikzpicture}
\captionof{figure}{Transformações dos dados e braços que leem cada camada.}\label{fig:dados}
\end{minipage}\par\smallskip


\subsection{Arquiteturas comparadas}

Foram comparadas \nArms{} arquiteturas (braços), ordenadas pela complexidade do que há entre a pergunta e o
modelo, da ausência de dados à consulta ao grafo com ferramentas (Tabela~\ref{tab:bracos}). Todas usam o mesmo
modelo (\modelName{}, temperatura zero), o mesmo prompt de sistema, que não nomeia a partida, e a mesma saída
estruturada (jogadores, valor, indicação de ausência de dados e justificativa).

\vspace{2mm}
\noindent\begin{minipage}{\columnwidth}
\centering\small\renewcommand{\arraystretch}{1.15}
\begin{tabular}{>{\raggedright\arraybackslash}p{0.32\columnwidth}>{\raggedright\arraybackslash}p{0.56\columnwidth}} % lint: ignore
\toprule
Braço & O que o modelo recebe e quem calcula \\
\midrule
\armNoContext{} & só a pergunta; ninguém calcula \\
\armEventsInPrompt{} & a tabela de \nActions{} ações inteira; o modelo calcula \\
\armStatsInPrompt{} & totais por jogador e por time; já estão calculados \\
\armVector{} & as \nTopK{} linhas de eventos mais parecidas com a pergunta; o modelo calcula \\
\armTextToCypher{} & o esquema do grafo; o modelo escreve a consulta e a ferramenta a executa \\
\armGraphTools{} & \nTools{} ferramentas, cada uma com uma consulta fixa; o modelo escolhe a ferramenta e os parâmetros, e ela executa a consulta \\
\bottomrule
\end{tabular}
\captionof{table}{As \nArms{} arquiteturas, da menos à mais complexa.}\label{tab:bracos}
\end{minipage}\par\smallskip

O braço \armVector{} segue o padrão de recuperação aumentada por geração~\cite{lewis2020rag}. Em
\armTextToCypher{} e \armGraphTools{}, o modelo atua como agente: em vez de responder de uma vez, decide a cada
passo qual ação tomar, recebe o resultado e decide de novo, até responder~\cite{wang2024agentssurvey,yao2023react}.
A ação é a chamada de uma ferramenta, isto é, de uma função externa cujo nome e argumentos o modelo escolhe e
cujo resultado volta ao seu contexto~\cite{schick2023toolformer} (Figura~\ref{fig:agente}). Cada pergunta admite no
máximo \nMaxCalls{} chamadas.

% F2: o ciclo do agente nos braços de consulta e de ferramentas. Diagrama em TikZ.
\vspace{2mm}
\noindent\begin{minipage}{\columnwidth}
\centering\small
\begin{tikzpicture}[
  box/.style={draw, rounded corners=2pt, align=center, inner sep=3pt, minimum height=7mm},
  arr/.style={-{Stealth[length=2mm]}}]
\node[box, text width=0.24\columnwidth] (q) at (-0.36\columnwidth,0) {Pergunta};
\node[box, text width=0.24\columnwidth] (m) at (0,0) {Modelo};
\node[box, text width=0.24\columnwidth] (a) at (0.36\columnwidth,0) {Resposta};
\node[box, text width=0.5\columnwidth] (t) at (0,-1.8) {Ferramenta};
\node[box, text width=0.5\columnwidth] (g) at (0,-3.4) {Grafo Neo4j};
\draw[arr] (q) -- (m);
\draw[arr] (m) -- (a);
\draw[arr] ([xshift=-4mm]m.south) -- node[left] {chamada} ([xshift=-4mm]t.north);
\draw[arr] ([xshift=4mm]t.north) -- node[right] {resultado} ([xshift=4mm]m.south);
\draw[arr] ([xshift=-4mm]t.south) -- ([xshift=-4mm]g.north);
\draw[arr] ([xshift=4mm]g.north) -- ([xshift=4mm]t.south);
\node[align=left] at (0.385\columnwidth,-1.8) {até \nMaxCalls{}\\ chamadas};
\end{tikzpicture}
\captionof{figure}{Ciclo do agente em \armTextToCypher{} e \armGraphTools{}. No primeiro, a ferramenta executa no grafo a consulta Cypher que o modelo escreve; no segundo, o modelo escolhe uma das \nTools{} ferramentas e seus parâmetros, e ela executa uma consulta fixa.}\label{fig:agente}
\end{minipage}\par\smallskip


Os dois braços de agente consultam o mesmo grafo e diferem em de onde vem a consulta. Em \armTextToCypher{}, a
ferramenta executa no grafo a consulta Cypher que o modelo escreve~\cite{ozsoy2025text2cypher}, sem a biblioteca de
algoritmos de grafos (GDS)~\cite{neo4j_gds} e sem os padrões táticos já calculados; ela só aceita leitura e devolve
até \nMaxRows{} linhas. Em \armGraphTools{}, cada ferramenta executa no grafo uma consulta fixa, definida no código
da ferramenta; o modelo escolhe a ferramenta e preenche seus parâmetros (por exemplo, o time e o período de uma
rede de passes e a métrica a calcular). É a ideia de oferecer ao modelo uma definição de ferramenta para consultar
o banco, com filtros e agregações como argumentos~\cite{shorten2025dbfunction}.

\subsection{Perguntas e gabarito}

As \nQuestions{} perguntas, em linguagem de futebol e de resposta fechada (quem, quantos, quais), formam cinco
grupos (Tabela~\ref{tab:perguntas}). O grupo de controle reúne perguntas cuja resposta não existe nos dados, como
as que pedem velocidade ou distância percorrida; a resposta certa é a ausência de dados. Os resultados são
reportados com e sem ele, isto é, nas \nAnswerable{} perguntas com resposta e nas \nQuestions{}.

\vspace{2mm}
\noindent\begin{minipage}{\columnwidth}
\centering\small\renewcommand{\arraystretch}{1.15}
\begin{tabular}{>{\raggedright\arraybackslash}p{0.28\columnwidth}>{\raggedright\arraybackslash}p{0.60\columnwidth}} % lint: ignore
\toprule
Grupo & Exemplo (resposta esperada) \\
\midrule
Busca e contagem (\nGroupSearch{}) & Quem recebeu o primeiro cartão amarelo da final? (jogador) \\
Relações na rede de passes (\nGroupRelations{}) & Qual dupla da França mais trocou passes entre si na final? (conjunto de jogadores) \\
Jogadas e sequências (\nGroupPlays{}) & Qual jogador participou de mais jogadas da Argentina que terminaram em finalização? (jogador) \\
Antes e depois de um momento (\nGroupMoments{}) & Depois que Dembélé foi substituído no primeiro tempo, qual companheiro mais passou a bola para Mbappé? (jogador) \\
Controle (\nGroupControl{}) & Qual foi a distância total percorrida por Rodrigo De Paul na final? (ausência de dados) \\
\bottomrule
\end{tabular}
\captionof{table}{Os cinco grupos de perguntas, com o número de perguntas de cada um e um exemplo.}\label{tab:perguntas}
\end{minipage}\par\smallskip

O gabarito foi extraído dos dados brutos. As perguntas foram escritas manualmente e suas respostas, validadas sobre
os dados brutos; foram eliminadas as perguntas com maior chance de ser interpretadas de mais de uma maneira.

\subsection{Saída estruturada e conferência}

Todos os braços respondem no mesmo formato: uma justificativa, uma lista de jogadores, um valor numérico e um
indicador de ausência de dados. A justificativa é escrita antes da resposta e nunca é pontuada. Uma resposta vazia
(sem jogador, sem valor e sem o indicador) é inválida e o modelo é reperguntado até duas vezes; se a resposta
continuar inválida, a execução conta como erro de formato, e não como resposta errada.

A conferência com o gabarito é feita por código, sem modelo como juiz, e depende do tipo da pergunta
(Tabela~\ref{tab:conferencia}). O formato pede que os rankings sejam listados do melhor ao pior; por isso, nas
perguntas que esperam um jogador, vale o primeiro jogador da lista e os demais não são conferidos. Os nomes são
normalizados (caixa e acentos) e aceitos pelo nome da escalação, pelo apelido ou pelo sobrenome, este só quando é
único entre os jogadores da partida.

\vspace{2mm}
\noindent\begin{minipage}{\columnwidth}
\centering\small\renewcommand{\arraystretch}{1.15}
\begin{tabular}{>{\raggedright\arraybackslash}p{0.26\columnwidth}>{\raggedright\arraybackslash}p{0.62\columnwidth}} % lint: ignore
\toprule
Tipo (perguntas) & Regra de acerto \\
\midrule
Jogador (\nCheckPlayer{}) & o primeiro jogador da lista é o do gabarito \\
Conjunto (\nCheckSet{}) & a lista tem exatamente os jogadores do gabarito, em qualquer ordem \\
Valor (\nCheckValue{}) & o número é igual ao do gabarito; quando a pergunta pede também um jogador, ele precisa conferir \\
Ausência de dados (\nCheckNoData{}) & o modelo indica que não há dados \\
\bottomrule
\end{tabular}
\captionof{table}{Os tipos de conferência, com o número de perguntas de cada um.}\label{tab:conferencia}
\end{minipage}\par\smallskip

\subsection{Resultados de uma execução e análise estatística}

Cada execução (uma pergunta feita a um braço) termina em um de quatro resultados (Tabela~\ref{tab:resultados}).
O termo alucinação é usado aqui em sentido operacional, mais estreito que o da literatura~\cite{ji2023survey}:
uma resposta errada apresentada como resposta. A abstenção é o oposto prudente e igualmente errado: o modelo
declara que não há dados em uma pergunta que tem resposta. Além do resultado, registram-se tokens, chamadas de
ferramenta e latência.

\vspace{2mm}
\noindent\begin{minipage}{\columnwidth}
\centering\small\renewcommand{\arraystretch}{1.15}
\begin{tabular}{>{\raggedright\arraybackslash}p{0.26\columnwidth}>{\raggedright\arraybackslash}p{0.62\columnwidth}} % lint: ignore
\toprule
Resultado & Definição \\
\midrule
Acerto & a resposta confere com o gabarito; no controle, o modelo indica que não há dados \\
Alucinação & o modelo dá uma resposta e ela está errada, inclusive um valor para uma pergunta do controle \\
Abstenção & o modelo indica que não há dados em uma pergunta que tem resposta \\
Erro de formato & não há resposta estruturada válida ao fim das tentativas \\
\bottomrule
\end{tabular}
\captionof{table}{Os quatro resultados possíveis de uma execução.}\label{tab:resultados}
\end{minipage}\par\smallskip

% TODO autor: REVISAR este parágrafo (qual o papel da análise estatística na Metodologia? talvez mover para Resultados)
\begingroup\color{red}\textbf{[REVISAR]} O acerto de um braço é a proporção de perguntas respondidas corretamente. Como as perguntas são poucas e cada uma
foi executada uma vez, duas análises acompanham essa proporção, ambas calculadas depois da rodada, a partir do
arquivo de resultados. A primeira é o intervalo de confiança de Wilson~\cite{wilson1927probable}, de \ciLevel{}
para cada acerto; diferentemente do intervalo normal, ele não ultrapassa zero nem cem por cento, o que importa para
os braços em zero por cento. A segunda compara pares de braços que responderam às mesmas perguntas: o teste exato
de McNemar~\cite{mcnemar1947note} considera apenas as perguntas em que um braço acertou e o outro errou, e é o que
Dietterich~\cite{dietterich1998approximate} indica quando o algoritmo é executado uma só vez. Foram comparados
\armTextToCypher{} e \armEventsInPrompt{}, \armGraphTools{} e \armTextToCypher{}, e \armGraphTools{} e
\armEventsInPrompt{}. Como são três testes, os valores $p$ recebem o ajuste de Holm~\cite{holm1979simple}, que
mantém no nível escolhido a probabilidade de algum falso positivo no conjunto. A rodada tem \nRepeats{} repetição
por pergunta; as conclusões são lidas como indícios, e não como demonstração.\par\endgroup

\subsection{Escopo}

Ficam fora do escopo: várias partidas, dados de rastreamento ou fisiológicos, inferência sobre lesão, fadiga ou
saúde, valor de mercado, escalação, predição de resultado, qualidade da explicação tática, comparação entre modelos
e agentes que escrevem código. Os dados são abertos, de uma partida pública, e não incluem dado pessoal sensível;
% TODO autor: confirmar o parágrafo de LGPD (o rascunho do colega não está no repositório)
o trabalho não infere condição de saúde de atleta.

% Seção: Experimentos e resultados  (orçamento: 2.5 página(s))
%
% Narrativa combinada com o autor (13/10/2026): Ferramentas mantêm o acerto nas duas faixas de preço; Consulta e
% Contexto longo dependem de modelos de preço alto, com raciocínio, e pagam em tokens de saída e latência; tool calling
% efetivo é condição. Os Resultados descrevem; a Discussão interpreta. Sem McNemar no texto por ora (a tabela
% generated/tables/mcnemar_models.tex continua gerada). Sem preço em dólar (os preços mudam): só tokens e latência.

\section{Experimentos e resultados}\label{sec:experimentos}

\subsection{Descrição dos dados}

Os dados são os eventos da final da Copa do Mundo de 2022: \nEvents{} eventos no JSON bruto, que viraram
\nActions{} ações na tabela de ações. As amostras do experimento são as perguntas: \nQuestions{} perguntas em cinco
grupos (Tabela~\ref{tab:perguntas}), das quais \nAnswerable{} têm resposta nos dados e \nControl{} formam o grupo de
controle. Cada pergunta foi feita a cada uma das \nArms{} abordagens, num total de \nRuns{} execuções por modelo.

\vspace{2mm}
\noindent\begin{minipage}{\columnwidth}
\centering\small\renewcommand{\arraystretch}{1.15}
\begin{tabular}{>{\raggedright\arraybackslash}p{0.28\columnwidth}>{\raggedright\arraybackslash}p{0.60\columnwidth}} % lint: ignore
\toprule
Grupo & Exemplo (resposta esperada) \\
\midrule
Busca e contagem (\nGroupSearch{}) & Quem recebeu o primeiro cartão amarelo da final? (jogador) \\
Relações na rede de passes (\nGroupRelations{}) & Qual dupla da França mais trocou passes entre si na final? (conjunto de jogadores) \\
Jogadas e sequências (\nGroupPlays{}) & Qual jogador participou de mais jogadas da Argentina que terminaram em finalização? (jogador) \\
Antes e depois de um momento (\nGroupMoments{}) & Depois que Dembélé foi substituído no primeiro tempo, qual companheiro mais passou a bola para Mbappé? (jogador) \\
Controle (\nGroupControl{}) & Qual foi a distância total percorrida por Rodrigo De Paul na final? (ausência de dados) \\
\bottomrule
\end{tabular}
\captionof{table}{Os cinco grupos de perguntas, com o número de perguntas de cada um e um exemplo.}\label{tab:perguntas}
\end{minipage}\par\smallskip

\subsection{\textit{Setup} experimental}

Foram avaliados \nModels{} modelos de linguagem, em duas faixas de preço de tabela (outubro de 2026). Na faixa
baixa, de centavos de dólar por milhão de tokens, ficaram \modelHaiku{} e \modelQwen{}; na alta, de alguns dólares
por milhão, \modelSonnet{} e \modelKimi{}. Em cada faixa há um modelo de código fechado e um de pesos abertos. Cada
modelo respondeu a todas as perguntas em todas as abordagens, com os mesmos prompts. Os modelos de preço baixo rodaram
sem raciocínio estendido e com temperatura zero; os de preço alto, com raciocínio estendido no nível padrão e na
temperatura padrão do provedor. A resposta de cada chamada ao modelo, raciocínio incluído, foi limitada a
\nMaxOutTokens{} tokens. Na abordagem \armEventsInPrompt{} do \modelSonnet{}, o raciocínio no nível padrão esgotou
esse limite em parte das perguntas, que ficaram sem resposta; nessa abordagem, o nível de raciocínio foi reduzido, e
as perguntas que ainda esgotaram o limite foram refeitas com o dobro dele.

A abordagem \armVector{} usou embeddings do modelo \texttt{\embedModel}. Nas abordagens com agente, cada pergunta
admitiu no máximo \nMaxCalls{} chamadas de ferramenta. A rodada de cada modelo teve \nRepeats{} repetição por
pergunta.

% TODO autor: REVISAR este parágrafo (movido da Metodologia; ainda não decidido se entra no artigo)
\begingroup\color{red}\textbf{[REVISAR]} O acerto de uma abordagem é a proporção de perguntas respondidas corretamente. Como as perguntas são poucas e cada uma
foi executada uma vez, duas análises acompanham essa proporção, ambas calculadas depois da rodada, a partir do
arquivo de resultados. A primeira é o intervalo de confiança de Wilson~\cite{wilson1927probable}, de \ciLevel{}
para cada acerto; diferentemente do intervalo normal, ele não ultrapassa zero nem cem por cento, o que importa para
as abordagens em zero por cento. A segunda compara pares de abordagens que responderam às mesmas perguntas: o teste exato
de McNemar~\cite{mcnemar1947note} considera apenas as perguntas em que uma abordagem acertou e a outra errou, e é o que
Dietterich~\cite{dietterich1998approximate} indica quando o algoritmo é executado uma só vez. Foram comparados
\armTextToCypher{} e \armEventsInPrompt{}, \armGraphTools{} e \armTextToCypher{}, e \armGraphTools{} e
\armEventsInPrompt{}. Como são três testes, os valores $p$ recebem o ajuste de Holm~\cite{holm1979simple}, que
mantém no nível escolhido a probabilidade de algum falso positivo no conjunto. A rodada tem \nRepeats{} repetição
por pergunta; as conclusões são lidas como indícios, e não como demonstração.\par\endgroup

% TODO autor: REVISAR (confirmar se o repositório é público e o link; o template pede aqui as etapas para reproduzir)
\begingroup\color{red}\textbf{[REVISAR]} O código, as perguntas, o gabarito, os resultados de cada execução e uma
página por rodada estão em \url{https://github.com/joaowinderfeldbussolotto/football-tactics-knowledge-graph}. As
bibliotecas usadas foram socceraction~\cite{socceraction}, Neo4j Graph Data Science~\cite{neo4j_gds},
networkx~\cite{networkx}, PydanticAI~\cite{pydanticai} e Langfuse~\cite{langfuse}, nas versões indicadas nas
referências.\par\endgroup

\subsection{Análise dos resultados}

\subsubsection{Acerto por abordagem e modelo}

A Tabela~\ref{tab:acerto} mostra o acerto das \nArms{} abordagens nos \nModels{} modelos, e a
Figura~\ref{fig:modelos} destaca as três que leem os dados da partida. \armGraphTools{} acertou de
\accRangeGraphTools{} das perguntas com resposta, nas duas faixas de preço. \armTextToCypher{} acertou de
\accRangeLowTextToCypher{} nos modelos de preço baixo e de \accRangeHighTextToCypher{} nos de preço alto;
\armEventsInPrompt{}, de \accRangeLowEventsInPrompt{} e de \accRangeHighEventsInPrompt{}. Nos modelos de preço
baixo, \armGraphTools{} ficou à frente das outras duas, com intervalos que não se sobrepõem aos de
\armEventsInPrompt{}; nos de preço alto, os intervalos das três se sobrepõem. As outras abordagens ficaram
abaixo em todos os modelos: \armVector{} acertou de \accRangeVector{}, \armStatsInPrompt{} de
\accRangeStatsInPrompt{} e \armNoContext{}, \accRangeNoContext{}.

\vspace{2mm}
\noindent\begin{minipage}{\columnwidth}
\centering\footnotesize\renewcommand{\arraystretch}{1.15}\setlength{\tabcolsep}{4pt}
\gentable{models}
\captionof{table}{Acerto, em porcentagem, nas perguntas com resposta. Os intervalos de confiança das três
abordagens que leem os dados estão na Figura~\ref{fig:modelos}.}\label{tab:acerto}
\end{minipage}\par\smallskip

\vspace{2mm}
\noindent\begin{minipage}{\columnwidth}
\centering
\gentable{fig_models}
\captionof{figure}{Acerto das três abordagens que leem os dados nas perguntas com resposta, por modelo, com o
intervalo de Wilson de \ciLevel{}. A área cinza marca os modelos de preço alto.}\label{fig:modelos}
\end{minipage}\par\smallskip

\subsubsection{Acerto por grupo de perguntas}

A Tabela~\ref{tab:grupos} mostra os acertos das três abordagens em cada grupo de perguntas
(Seção~\ref{sec:perguntas}). Nos modelos de preço baixo, a distância entre \armGraphTools{} e as outras duas
abordagens aparece sobretudo nos grupos de jogadas e de relações na rede de passes. Nos modelos de preço alto,
\armGraphTools{} acertou menos que \armTextToCypher{} no grupo antes e depois de um momento do jogo.

\vspace{2mm}
\noindent\begin{minipage}{\columnwidth}
\centering\footnotesize\renewcommand{\arraystretch}{1.15}\setlength{\tabcolsep}{4pt}
\gentable{groups_models}
\captionof{table}{Acertos por grupo de perguntas, em cada célula na ordem \armShortEventsInPrompt{},
\armShortTextToCypher{} e \armShortGraphTools{}. Busca: busca e contagem; Relações: relações na rede de passes;
Jogadas: jogadas e sequências; Momento: antes e depois de um momento do jogo.}\label{tab:grupos}
\end{minipage}\par\smallskip

\subsubsection{Erros e custo}

A Tabela~\ref{tab:erros} mostra, para as três abordagens que leem os dados, a proporção de alucinações, de
abstenções e de erros de formato (Seção~\ref{sec:conferencia}) e, por execução, a média de tokens de entrada e de
saída e a latência. Nos modelos de preço alto, \armEventsInPrompt{} produziu em média de
\outRangeHighEventsInPrompt{} tokens de saída por pergunta, quase todos de raciocínio, contra
\outRangeLowEventsInPrompt{} nos de preço baixo, e levou de \latRangeHighEventsInPrompt{} segundos. Os erros de
formato vieram de limites: no \modelQwen{}, o agente esgotou as chamadas de ferramenta em \armTextToCypher{}
(\fmtQwenTextToCypher{}) e em \armGraphTools{} (\fmtQwenGraphTools{}); no \modelKimi{}, em \armGraphTools{}
(\fmtKimiGraphTools{}), o modelo esgotou o limite de saída ou as tentativas de resposta; no \modelSonnet{}, em
\armEventsInPrompt{} (\fmtSonnetEventsInPrompt{}), esgotou o limite de saída mesmo com o dobro dele. Fora da tabela, \armNoContext{} e \armStatsInPrompt{} se abstiveram em
\abstRangeNoContext{} e \abstRangeStatsInPrompt{} das perguntas, e \armVector{}, em \abstRangeVector{}.

\vspace{2mm}
\noindent\begin{minipage}{\columnwidth}
\centering\footnotesize\renewcommand{\arraystretch}{1.15}\setlength{\tabcolsep}{3pt}
\gentable{errors_models}
\captionof{table}{Perfil de erro sobre todas as perguntas e médias por execução das três abordagens que leem os
dados, por modelo. Form.: erro de formato; entrada e saída em milhares de tokens; latência em segundos.}\label{tab:erros}
\end{minipage}\par\smallskip

% Seção: Conclusão  (orçamento: 0,5 página)
%
% Responde às três perguntas de pesquisa da Introdução e aponta direções futuras ligadas às limitações da
% Discussão. Sem frase de fechamento.

\section{Conclusão}\label{sec:conclusão}

Este trabalho comparou \nArms{} abordagens de aprendizado em contexto para responder a perguntas sobre os eventos
da final da Copa do Mundo de 2022, em \nModels{} modelos. Nenhuma abordagem acertou mais em todos os modelos, mas a
abordagem \armGraphTools{} foi a única que manteve o acerto alto em todos, de \accRangeGraphTools{}. A \armTextToCypher{} e o
\armEventsInPrompt{} variaram entre os modelos e ficaram abaixo dela nos menores. Nos maiores, com raciocínio
estendido, a \armTextToCypher{} chegou ao mesmo nível, e o \armEventsInPrompt{} se aproximou, este ao custo de mais
tokens de saída e mais tempo por pergunta. A \armVector{} e as \armStatsInPrompt{} acertaram pouco em todos os modelos.

Os resultados vêm de uma partida e de uma repetição por pergunta. Repetir as execuções permitiria medir a variação
entre rodadas e testar as diferenças entre abordagens. Outros caminhos são estender as perguntas a várias partidas;
avaliar modelos maiores sem raciocínio estendido, para separar o efeito do tamanho do efeito do raciocínio; aplicar
as mesmas ferramentas a um banco relacional (SQL), sem o grafo, para separar o efeito das ferramentas do efeito do
grafo; e testar uma busca vetorial em que cada trecho traga também os eventos vizinhos, como em Tzikas et
al.~\cite{tzikas2025dynamic}. Por fim, as mesmas abordagens podem ser testadas em outros registros de eventos, como
os de processos industriais, com a preparação dos dados, o grafo e as ferramentas feitos para as perguntas desse
domínio.
% Declaração de uso de IA generativa (exigência da disciplina). Formato das políticas da Elsevier e da SciELO:
% ferramenta, finalidade, revisão pelos autores e responsabilidade. Logo antes das referências.

\section{Declaração de uso de IA generativa}\label{sec:declaração}

Durante a preparação deste trabalho, os autores usaram ferramentas de IA generativa da Anthropic e da OpenAI
como assistentes na implementação do código, na busca e conferência de referências e na redação e revisão do
texto. Os autores revisaram e editaram todo o conteúdo e assumem total responsabilidade por ele.


\bibliographystyle{plain}
\bibliography{refs}

\end{multicols*}

\end{document}
